\section{History-dependent interaction and revision rights}\label{sec:model}

\subsection{From a directional field to a history payoff}
Consider two formations, Blue and Red, moving in the plane over a finite horizon
$T$. Blue maximizes net interaction payoff and Red minimizes it. A leader's
motion is generated by a sequence of bounded turn controls. Trailing vessels
occupy fixed arc-length offsets along its executed track; their orientations
are the local track tangents. For formation $i$, vessel $k$ has position and
heading
\begin{equation}
 x_{ik}(t)=\gamma_i(s_i(t)-\ell_k),\qquad
 \psi_{ik}(t)=\arg\gamma_i'(s_i(t)-\ell_k).
 \label{eq:formation}
\end{equation}
The numerical implementation uses a causal polyline and its incoming tangent
at an exact station. A straight prehistory supplies the trailing offsets at
the initial epoch. Current leader pose alone is generally insufficient to
recover the trailing configuration. The analysis therefore retains executed
histories rather than identifying states solely by leader position and heading.

Let $r_{kl}$ denote pairwise separation, and $\delta^B_{kl},\delta^R_{lk}$
the target bearings measured relative to the respective vessel headings.
A directional kernel $K_i$ gives the interaction intensity for each pair.
The formation-level differential is
\begin{equation}
 L(h_t)=\sum_{k,l}\left[
 K_B(r_{kl},\delta^B_{kl},\delta^R_{lk})-
 K_R(r_{kl},\delta^R_{lk},\delta^B_{kl})\right].
 \label{eq:field}
\end{equation}
Here $h_t$ is the complete executed history through time $t$. The specific
kernel used below is a frozen surrogate of an audited game-rule response,
not a model calibrated to observed naval engagements. Its firing sectors
are piecewise defined; no global differentiability is assumed.

The objective accumulates interaction along the trajectory. The finite
benchmark uses the causal endpoint trapezoid
\begin{equation}
 J(h_T)=\frac12\sum_{t=0}^{T-1}\{L(h_t)+L(h_{t+1})\}.
 \label{eq:payoff}
\end{equation}
The initial term is a common constant for a fixed instance; every command
affects a subsequent state included in the same evaluation horizon. This
is a specified discrete objective, not an assertion of exact continuous-time
integration. Directional motion can change the ordering of paths because
an attractive terminal configuration need not compensate for unfavorable
interaction accumulated earlier. We compare the integrated and terminal
objectives under the same horizon in a separate ablation.

Let each player have a state-independent finite action menu of size $b$.
Enumerating the $b^T$ complete action paths of each player gives a matrix
$A\in\mathbb R^{b^T\times b^T}$, where $A(i,j)=J(h_T(i,j))$.
This matrix indexes realized physical paths. Its rows are \emph{not} all
feedback policies: a policy may select different own paths after different
observed opponent prefixes. The distinction is essential when computing
adaptation value.

\subsection{Public calendars and sealed command blocks}
Write $\NN=\{1,\ldots,T-1\}$. A calendar $S\subseteq\NN$ specifies Blue's
revision opportunities; an initial decision at $t=0$ is compulsory and does
not consume the revision budget. Calendars are fixed and publicly known
before play. At each $t\in\{0\}\cup S$, Blue observes both executed action
prefixes and selects a complete block of commands until its next update.
Red follows an independently specified calendar $R$. The simultaneous current
commands and the unexecuted portions of either block are hidden from the
opponent. Between updates a player executes its chosen block. At its next
update it remembers its earlier choices and the complete executed history.

This is a finite perfect-recall extensive-form game. Revision rights constrain
when commands may respond to observations; they do not erase a player's
memory. No noisy observations, damage states, private capability types or
history-dependent legal-action restrictions are added to the numerical model.

Let $\PP_S$ and $\QQ_R$ denote the mixed-policy sets induced by the calendars.
Define
\begin{equation}
 V(S,R)=\max_{p\in\PP_S}\min_{q\in\QQ_R}\EE_{p,q}[A(i,j)].
 \label{eq:value}
\end{equation}
Finite minimax supplies a value. We study two fixed opponent conditions:
$R=\varnothing$ (committed, denoted $C$) and $R=\NN$ (flexible, denoted $F$).
Results are reported separately. In a calendar comparison, the physical
matrix and the opponent's policy space remain unchanged.

\subsection{The budget frontier and the decision target}
For $K\in\{0,\ldots,T-1\}$, the calendar design problem is
\begin{equation}
 W_K(R)=\max_{S\subseteq\NN:\,|S|\le K}V(S,R).
 \label{eq:budget}
\end{equation}
The design variable is a deterministic public calendar. The opponent may
respond strategically within its fixed policy class. Hidden randomized
calendars and jointly optimized opposing calendars define different design
problems and are outside the present formulation.

Let $V_0=V(\varnothing,R)$, $V_F=V(\NN,R)$ and $D=V_F-V_0$.
When $D$ is numerically resolved and positive, the smallest budget retaining
a fraction $\rho\in(0,1)$ of full adaptation value is
\begin{equation}
 K_\rho=\min\{K:W_K(R)\ge\rho V_F+(1-\rho)V_0\}.
 \label{eq:retention}
\end{equation}
Otherwise we report absolute performance loss. This prevents an almost-zero
adaptation premium from producing unstable percentage comparisons.

\begin{table}[t]
\centering\small
\caption{Main notation and assumptions. All calendar comparisons hold the physical instance and the opponent condition fixed.}\label{tab:notation}
\begin{tabular}{p{.18\linewidth}p{.73\linewidth}}
\toprule
Symbol & Meaning \\
\midrule
$T,b$ & Horizon and finite per-epoch action-menu size \\
$A(i,j)$ & Accumulated payoff for a complete pair of executed paths \\
$S,R$ & Public Blue and Red revision calendars; initial decision excluded \\
$\PP_S,\QQ_R$ & Perfect-recall policies with hidden unexecuted command blocks \\
$V(S,R),W_K(R)$ & Fixed-calendar value and best value with at most $K$ revisions \\
$[L_S,U_S]$ & Recorded finite-game numerical value interval \\
$K_\rho$ & Minimal budget retaining fraction $\rho$ of full adaptation \\
Standing assumptions & Known initial state, deterministic causal dynamics, bounded payoff, state-independent legal action menus, full executed-history observation at updates \\
\bottomrule
\end{tabular}
\end{table}
